\documentclass[11pt]{article}
\usepackage[localtheorems]{raeez-math-template}

\newtheorem{theorem}{Theorem}
\newtheorem{proposition}[theorem]{Proposition}
\newtheorem{conjecture}[theorem]{Conjecture}
\newtheorem{definition}[theorem]{Definition}
\newtheorem{remark}[theorem]{Remark}

\providecommand{\K}{\mathrm{K3}}
\providecommand{\KE}{\K\times E}
\providecommand{\C}{\mathbb{C}}
\providecommand{\Z}{\mathbb{Z}}
\providecommand{\HH}{\mathbb{H}}
\providecommand{\Sp}{\mathrm{Sp}}
\providecommand{\OO}{\mathrm{O}}
\providecommand{\Aut}{\mathrm{Aut}}
\providecommand{\Rep}{\mathrm{Rep}}
\providecommand{\CoHA}{\mathrm{CoHA}}
\providecommand{\DT}{\mathrm{DT}}
\providecommand{\red}{\mathrm{red}}
\providecommand{\wt}{\operatorname{wt}}
\providecommand{\Tr}{\operatorname{Tr}}
\providecommand{\Borch}{\operatorname{Borch}}
\providecommand{\mult}{\operatorname{mult}}
\providecommand{\kBKM}{\kappa_{\mathrm{BKM}}}
\providecommand{\kch}{\kappa_{\mathrm{ch}}}
\providecommand{\kcat}{\kappa_{\mathrm{cat}}}
\providecommand{\kfib}{\kappa_{\mathrm{fiber}}}
\providecommand{\kHeis}{\kappa_{\mathrm{ch}}^{\mathrm{Heis}}}
\providecommand{\cO}{\mathcal{O}}
\providecommand{\cZ}{\mathcal{Z}}
\providecommand{\fg}{\mathfrak{g}}
\providecommand{\Mtf}{M_{24}}
\providecommand{\Dfive}{\Delta_5}
\providecommand{\Phiten}{\Phi_{10}}
\providecommand{\fone}{\phi_{0,1}}
\providecommand{\Zg}{Z_g}

\title{Mathieu Twining, CHL Denominators, and Lorgat Conjecture 1}
\author{Raeez Lorgat}
\date{}

\begin{document}
\maketitle

\begin{abstract}
Mathieu moonshine, twisted-twined K3 elliptic genera, CHL denominator
constants, Gritsenko and Clery diagonal-divisor products, Igusa square
normalisations, and reduced \(K3\times E\) Donaldson-Thomas series are
linked by precise comparison maps.  They are indexed by different
objects.  A row of Lorgat Conjecture 1 is meaningful only after the
Mathieu datum, Jacobi input, multiplier, Borcherds lift, denominator
lattice, Weyl chamber, and reduced DT scalar normalisation have all
been fixed.  In the primitive CHL denominator normalisation
\[
  c_N(0)=(10,8,6,4,2),\qquad
  \kBKM(\Phi_N)=c_N(0)/2=(5,4,3,2,1),
  \quad N\in\{1,2,3,4,6\}.
\]
\end{abstract}

\tableofcontents

\section{Objects}

\begin{definition}[Ordinary Mathieu twining]
For a conjugacy class \([g]\subset\Mtf\), the ordinary twined K3
elliptic genus is
\[
  \Zg(\tau,z)
  =
  \Tr_{\mathcal H_{RR}}
  \left(g(-1)^{F+\widetilde F}q^{L_0-c/24}
  \bar q^{\widetilde L_0-\widetilde c/24}y^{J_0}\right).
\]
It is a weak Jacobi form of weight \(0\), index \(1\), and level
dividing \(|g|\), with multiplier determined by the Frame shape of
\(g\).  Its \(q^0\)-term is written
\[
  \Zg(\tau,z)=2y^{-1}+c^{(g)}(0,0)+2y+O(q).
\]
\end{definition}

\begin{definition}[Twisted-twined genus]
A twisted-twined genus \(\phi_{g,h}\) is defined for a commuting pair
\(gh=hg\).  It is the trace of \(h\) on the \(g\)-twisted sector.
The Persson-Volpato anomaly is the transgression of a class in
\(H^3(\Mtf,U(1))\) to the centraliser of \(g\).  The pair
\((g,h)\), the centraliser class of \(h\), and the anomaly multiplier
are part of the datum.
\end{definition}

\begin{definition}[Borcherds input]
A Borcherds denominator input is a tuple
\[
  (F,\Gamma,\nu,L,\rho,\Delta^+,\mu),
\]
where \(F=\sum c(n,\ell)q^ny^\ell\) is the Jacobi or vector-valued
modular input, \(\Gamma\) is the modular group or cover, \(\nu\) is
the multiplier, \(L\) is the Lorentzian or orthogonal lattice, \(\rho\)
is the Weyl vector, \(\Delta^+\) is the product chamber, and
\(\mu\) is the signed root-multiplicity convention.  Borcherds'
singular theta lift gives
\[
  \Phi=\Borch(F),\qquad
  \kBKM(\Phi):=\wt(\Phi)=\frac{c(0)}2
\]
at the chosen cusp.
\end{definition}

\section{Mathieu Twining Table}

The ordinary \(M_{24}\) twining constants are class constants, separate
from CHL denominator constants.
\[
\begin{array}{c|rrrrrrrrrrrrr}
[g] &1A&2A&2B&3A&3B&4A&4B&4C&5A&6A&6B&7A&7B\\
\hline
c^{(g)}(0,0)&20&4&-4&2&-4&0&4&-4&0&0&0&-1&-1
\end{array}
\]
\[
\begin{array}{c|rrrrrrrrrrrrr}
[g] &8A&10A&11A&12A&12B&14A&14B&15A&15B&21A&21B&23A&23B\\
\hline
c^{(g)}(0,0)&0&-4&-2&0&0&0&0&0&0&-1&-1&0&0
\end{array}
\]
The anomalous ordinary classes for twisted-twined purposes are
\[
  7A,\ 7B,\ 11A,\ 23A,\ 23B.
\]
For these classes the ordinary twining genus exists, while a
twisted-twined denominator row requires the transgressed projective
multiplier on the selected centraliser sector.

\section{Primitive CHL Denominators}

The primitive CHL denominator branch used in the compact
\(K3\times E\) Borcherds lane is
\[
  N\in\{1,2,3,4,6\}.
\]
Its constants and weights are
\[
\begin{array}{c|c|c|c}
N & \text{Frame shape} & c_N(0) & \kBKM(\Phi_N)\\
\hline
1 & 1^{24} & 10 & 5\\
2 & 1^8 2^8 & 8 & 4\\
3 & 1^6 3^6 & 6 & 3\\
4 & 1^4 2^2 4^4 & 4 & 2\\
6 & 1^2 2^2 3^2 6^2 & 2 & 1
\end{array}
\]
Singly twined scalar lists on the same order set arise from different
input choices and belong outside the primitive CHL denominator
normalisation.

\begin{proposition}[CHL Borcherds weight]
For \(N\in\{1,2,3,4,6\}\),
\[
  \kBKM(\Phi_N)=\wt(\Phi_N)=\frac{c_N(0)}2,
\]
and hence
\[
  (\kBKM(\Phi_1),\kBKM(\Phi_2),\kBKM(\Phi_3),
  \kBKM(\Phi_4),\kBKM(\Phi_6))=(5,4,3,2,1).
\]
\end{proposition}

\begin{proof}
Borcherds' product theorem gives weight \(c(0)/2\) for the singular
theta lift of the chosen input.  The CHL-averaged inputs have
constant terms \(10,8,6,4,2\) at \(N=1,2,3,4,6\).  Substitution gives
the displayed weights.  The computation uses the CHL-averaged Jacobi
input, separate from the ordinary \(c^{(g)}(0,0)\) row of the
\(M_{24}\) table.
\end{proof}

\section{Diagonal-Divisor Forms}

The Gritsenko and Clery diagonal-divisor catalogue is indexed by
triples \((t,N;k)\):
\[
\begin{array}{c|c|c|c}
\text{form} & (t,N) & \text{group} & \wt\\
\hline
\Dfive & (1,1) & \Sp_4(\Z) & 5\\
\Delta_2 & (2,1) & \Gamma_2 & 2\\
\nabla_3 & (1,2) & \Gamma_0^{(2)}(2) & 3\\
\Delta_1 & (3,1) & \Gamma_3 & 1\\
\nabla_2 & (1,3) & \Gamma_0^{(2)}(3) & 2\\
\Delta_{1/2} & (4,1) & \Gamma_4 & 1/2\\
\nabla_{3/2} & (1,4) & \Gamma_0^{(2)}(4) & 3/2\\
Q_1 & (2,2) & \Gamma_2(2) & 1
\end{array}
\]
The corresponding twice-weight constants are
\[
  (10,4,6,2,4,1,3,2).
\]
The common entry with the primitive CHL branch is the level-one
\(\Dfive\) row.

\section{Level One}

At \([g]=1A\), the full K3 elliptic genus is
\[
  Z_{1A}(\tau,z)=2\fone(\tau,z),
  \qquad
  \fone(\tau,z)=y^{-1}+10+y+O(q).
\]
The Borcherds input is the index-one generator \(\fone\), so
\[
  c_1(0)=10,\qquad
  \kBKM(\Dfive)=\frac{c_1(0)}2=5.
\]
The monic product and the theta-normalised Igusa square root satisfy
\[
  D_5=64^{-1}\Dfive,\qquad
  [q^{1/2}y^{1/2}p^{1/2}]\Dfive=64.
\]
The denominator algebra is
\[
  \operatorname{den}(\fg_{\Dfive})=D_5(2Z).
\]
The square normalisation is
\[
  \Phiten^{\mathrm{un}}=\Dfive^2,\qquad
  \Phiten^{\mathrm{OP}}=D_5^2=4096^{-1}\Dfive^2.
\]
The Oberdieck-Pandharipande reduced scalar convention for
\(K3\times E\) is
\[
  Z_{\mathrm{OP}/\DT}^{K3\times E}=-D_5^{-2}
  =-4096\,\Dfive^{-2}.
\]
Thus the primitive denominator is \(D_5\), while the scalar reduced
DT partition function uses its square inverse.

\section{Lorgat Conjecture 1}

\begin{conjecture}[Scoped row form]
Let \(F_j\) be one of the eight diagonal-divisor forms above.  A row
realisation of Lorgat Conjecture 1 consists of:
\begin{enumerate}[label=(\roman*)]
\item a compact Calabi-Yau or orbifold Calabi-Yau host \(X_j\) with
oriented reduced DT theory and fixed variables;
\item a scalar identity
\[
  Z_{\red,\DT}^{X_j}(Z)=C_jF_j(Z)^{-2}
\]
in a specified chamber;
\item a commuting Mathieu datum \((g,h)\) and its twisted-twined genus
\(\phi_{g,h}\);
\item a Borcherds input
\((\phi_{g,h},\Gamma_j,\nu_j,L_j,\rho_j,\Delta_j^+,\mu_j)\);
\item a denominator algebra \(\fg_j\) with
\[
  \operatorname{den}(\fg_j)=F_j,\qquad
  \mult_{\fg_j}(\alpha(n,\ell,m))
  =\mu_j\bigl([q^ny^\ell]\phi_{g,h}\bigr).
\]
\end{enumerate}
\end{conjecture}

The comparison map has two components.  The scalar component sends the
reduced DT series to \(F_j^{-2}\).  The denominator component sends
Jacobi coefficients of the selected twisted-twined genus to the signed
root multiplicities of \(\fg_j\).  A row assertion must name both
components.

\begin{proposition}[Level-one row]
The \(\Dfive\) row is supplied by
\[
  [g]=1A,\qquad (g,h)=(1A,1A),\qquad
  \phi_{g,h}=\fone,\qquad c_1(0)=10.
\]
Its Borcherds lift is \(D_5=64^{-1}\Dfive\), its BKM characteristic is
\(\kBKM(\Dfive)=5\), and its scalar \(K3\times E\) DT comparison uses
\(\Phiten^{\mathrm{OP}}=D_5^2\).
\end{proposition}

\begin{proof}
The \(1A\) twining genus is the ordinary K3 elliptic genus
\(2\fone\).  The primitive Borcherds input is \(\fone\), whose
discriminant-zero coefficient is \(10\).  Borcherds' weight theorem
gives weight \(5\).  Gritsenko-Nikulin identify the monic product with
\(D_5=64^{-1}\Dfive\).  Squaring gives the Igusa scalar convention
\(\Phiten^{\mathrm{OP}}=D_5^2\), and the reduced \(K3\times E\)
DT scalar is the reciprocal square with the Oberdieck-Pandharipande
sign.
\end{proof}

\section{The \(K3\times E\) Invariant Lanes}

For \(Y=K3\times E\), Kunneth gives
\[
  h^{0,\bullet}(Y)=(1,1,1,1),\qquad
  \kch(Y)=\sum_q(-1)^q h^{0,q}(Y)=0.
\]
The categorical invariant is
\[
  \kcat(Y)=\chi(\cO_{K3})\chi(\cO_E)=2\cdot 0=0.
\]
The remaining specialisations are
\[
  \kHeis(Y)=3,\qquad
  \kBKM(\fg_{\Dfive})=5,\qquad
  \kfib(K3)=24.
\]
These five values record five constructions:
\[
\begin{array}{c|c|c}
\text{construction} & \text{invariant} & \text{value}\\
\hline
\text{compact Hodge supertrace} & \kch(Y) & 0\\
\text{categorical Euler characteristic} & \kcat(Y) & 0\\
\text{Heisenberg-Mukai specialisation} & \kHeis(Y) & 3\\
\text{Borcherds denominator} & \kBKM(\fg_{\Dfive}) & 5\\
\text{K3 Mukai lattice rank} & \kfib(K3) & 24
\end{array}
\]
The rows of this table are independent constructions; no sum formula
relates them.

\section{Hall, Yangian, and Chiral Scope}

The BKM branch for \(K3\times E\) is the Hall-Drinfeld branch:
\[
  G(K3\times E)=D\bigl(Y^+_{\mathrm{Hall}}(K3\times E)\bigr)
\]
after the positive half, Hopf pairing, orientation data, and compact
Hall recognition gates have been constructed.  The Mukai Yangian
branch is the stable-envelope and Mukai-lattice branch attached to the
K3 surface.  The two branches have different inputs and different
comparison maps.

For complex dimension \(d\geq3\), the native specialised
CY-to-chiral output is \(E_1\)-chiral:
\[
  A_X=\mathrm{Sp}_{\Sigma_2,C}
  \bigl(\Phi^{\mathrm{FA}}_3(D^b\mathrm{Coh}(X))\bigr).
\]
The \(E_2\) structure belongs to the centre, Drinfeld double, or module
layer, for instance to
\[
  Z\bigl(\Rep^{E_1}(A_X)\bigr)
  \quad\text{or}\quad
  D\bigl(Y^+_{\mathrm{Hall}}(X)\bigr)
\]
when the corresponding hypotheses have been supplied.

For toric \(\C^3\), the positive-half statement is
\[
  \CoHA(\C^3)=Y^+(\widehat{\mathfrak{gl}}_1).
\]
The \(\mathcal W_{1+\infty}\) object is recovered through the
Drinfeld-centre or double passage.

\section{References}

\begin{itemize}
\item Borcherds 1998, Theorem 13.3: singular theta lift weight
\(c(0)/2\).
\item Gritsenko-Nikulin 1998 and Gritsenko 1999: \(\Dfive\),
the monic denominator product \(D_5\), and the rank-three BKM
denominator algebra.
\item Igusa, Maass, Freitag, and Van der Geer: \(\Phiten=\Dfive^2\),
\(\wt(\Dfive)=5\), and the Maass character of \(\Dfive\).
\item Eguchi-Ooguri-Tachikawa 2011 and Gaberdiel-Hohenegger-Volpato
2012: \(M_{24}\)-twined K3 elliptic genera.
\item Persson-Volpato 2012 and Gaberdiel-Persson-Volpato 2013:
twisted-twined genera and \(H^3(\Mtf,U(1))\) anomaly multipliers.
\item Oberdieck-Pandharipande and Oberdieck-Pixton: reduced
\(K3\times E\) DT scalar normalisation.
\item Clery and Gritsenko: the eight diagonal-divisor genus-two forms.
\end{itemize}

\end{document}
